% Source: https://artofproblemsolving.com/community/c3h2092838_amc_10_b_2004
% Stable identity: aops-topic-2092838-q01
Let the stack contain $p$ pennies, $n$ nickels, $d$ dimes and $q$ quarters. Then
\[1.55p+1.95n+1.35d+1.75q=14,\]
and multiplying by $20$ gives the integer equation
\[31p+39n+27d+35q=280 .\]
\emph{The step modulo $4$.} Each coefficient is one less than a multiple of $4$: $31=32-1$, $39=40-1$, $27=28-1$, $35=36-1$, so each is congruent to $-1$ modulo $4$; and $280=4\cdot70\equiv0$. Replacing every number by its remainder class modulo $4$ (which is allowed in sums and products) gives
\[-p-n-d-q\equiv0\pmod4,\]
that is, the number of coins $N=p+n+d+q$ is a multiple of $4$.

\emph{Finishing.} Every coin is between $1.35$ mm and $1.95$ mm thick, so $1.35N\le14\le1.95N$, which gives $7.2<N<10.4$. The only multiple of $4$ in this range is $N=8$, and $8$ coins are indeed possible (for example, $8$ quarters: $8\cdot1.75=14$). The answer is $\boxed{\textbf{(B)}\ 8}$.
